
        \documentclass[fleqn]{article}      
        \usepackage{amsmath}
        \usepackage{fontspec}
        \usepackage{kotex} % 한국어 지원  

        \begin{document}      
        쉬움: \\\\
1. 다음 중 적분의 정의에 가장 잘 맞는 것은?  
   A. 미분  
   B. 함수의 넓이  
   C. 기울기  
   D. 점의 위치 \\\\
   
2. 다음 함수의 적분을 구하시오: ∫x dx (x에 대한 적분) \\\\
   
3. 적분의 기호는 무엇인가요? \\\\
   
4. 다음 중 정적분에 대한 설명으로 올바른 것은?  
   A. 두 점 사이의 면적  
   B. 기울기를 계산하는 것  
   C. 함수의 극한  
   D. 점의 개수 \\\\
   
5. 적분의 결과는 항상 무엇으로 해석될 수 있나요?  
   A. 면적  
   B. 기울기  
   C. 점  
   D. 함수의 값 \\\\

중간: \\\\
1. 다음 함수 ∫(2x + 3) dx를 적분하시오. \\\\
   
2. 주어진 함수 f(x) = x²에 대하여 0부터 2까지의 정적분을 구하시오. \\\\
   
3. ∫(sin x) dx의 결과를 구하시오. \\\\
   
4. 어떤 함수의 정적분이 0이라면, 그 함수는 어떤 성질을 가지고 있나요? \\\\
   
5. 다음 중 부정적분의 예는 무엇인가요?  
   A. ∫(3x²) dx  
   B. ∫(1/x) dx  
   C. ∫(e^x) dx  
   D. 모두 정답 \\\\

어려움: \\\\
1. ∫(x³ - 4x + 1) dx의 결과를 구하시오. \\\\
   
2. 1부터 3까지의 구간에서 f(x) = x² - 2x의 정적분을 계산하시오. \\\\
   
3. ∫(e^(2x)) dx의 결과를 구하시오. \\\\
   
4. 다음 중 적분의 기본 정리에 대한 설명으로 올바른 것은?  
   A. 미분과 적분은 서로 반대이다.  
   B. 적분은 항상 양수이다.  
   C. 정적분은 부정적분의 결과와 같다.  
   D. 적분은 함수의 최대값을 찾는 것이다. \\\\
   
5. 어떤 함수 f(x)에 대해 ∫f'(x)dx = f(x) + C는 어떤 의미인가요? \\\\

      
        \end{document} 
    